% Faithful transcription — original German. Transcribed from the original printing:
% E. Noether, "Ableitung der Elementarteilertheorie aus der Gruppentheorie", Jahresbericht der
% Deutschen Mathematiker-Vereinigung 34 (1926), Abteilung 2 ("Mitteilungen und Nachrichten"),
% p. 104.
% Scan: GDZ (Göttingen), PPN37721857X_0034, scan image 361 = printed p. 104 of Abteilung 2.
%
% WHAT THIS ITEM IS, AND WHAT IS THEREFORE NOT TRANSCRIBED. The text is the abstract of a lecture
% Noether gave to the Mathematische Gesellschaft in Göttingen on 27 January 1925, printed as one
% run-in paragraph inside that society's collective meeting report. Printed p. 104 carries six
% such abstracts by six different speakers (Geiger, Ostrowski, Noether, Landau, Ore, Kneser);
% only Noether's is this work. The report's own apparatus is page furniture of the REPORT, not of
% this item, and is not transcribed: the running head, the page's four italic date headings
% (27. Januar 1925., 3. Februar 1925., 10. Februar 1925., 17. Februar 1925. — of which the
% first two bracket this item), and the em-dashes that separate one speaker's abstract from the
% next. The two em-dashes that ARE transcribed, around "z. B. Bettische und Torsionszahlen in
% der Topologie", are parenthetical dashes inside Noether's own final sentence.
%
% THE TITLE LINE. The print sets speaker and title run-in, as
%     — E. N o e t h e r: Ableitung der Elementarteilertheorie aus der Gruppentheorie. Die …
% with the surname letterspaced (Sperrung) and the title in plain roman, closed by a period, after
% which the abstract runs straight on within the same paragraph. House style gives the title a
% \section* and the attribution a \subsection*. The letterspacing is not carried into either:
% R25 forbids a brace group inside a heading, and R20 rules that a heading's own face carries its
% prominence. The period closing the title is the print's own and is kept (R12).
%
% \origpage{104} is the printed page number of ABTEILUNG 2, which is separately paginated from
% Abteilung 1 within the same bound volume — the volume's two runs are pp. 1–248 and pp. 1–184.
%
% No figures, no footnotes, no displayed equations. The item's only mathematics is the inline
% normal basis $(e_1 y_1, e_2 y_2, \ldots, e_r y_r)$, whose subscripts were read at magnification.
%
% See notation.md for the decisions a reviewer would otherwise have to re-derive from the scan.
\documentclass{article}
\usepackage{readmasters}
\begin{document}

\origpage{104}

\section*{Ableitung der Elementarteilertheorie aus der Gruppentheorie.}

\subsection*{E. Noether.}

Die Elementarteilertheorie gibt bekanntlich für Moduln aus ganzzahligen Linearformen eine Normalbasis von der Form $(e_1 y_1, e_2 y_2, \ldots, e_r y_r)$, wo jedes $e$ durch das folgende teilbar ist; die $e$ sind dadurch bis aufs Vorzeichen eindeutig festgelegt. Da jede Abelsche Gruppe mit endlich vielen Erzeugenden dem Restklassensystem nach einem solchen Modul isomorph ist, ist dadurch der Zerlegungssatz dieser Gruppen als direkte Summe größter zyklischer mitbewiesen. Es wird nun umgekehrt der Zerlegungssatz rein gruppentheoretisch direkt gewonnen, in Verallgemeinerung des für endliche Gruppen üblichen Beweises, und daraus durch Übergang vom Restklassensystem zum Modul selbst die Elementarteilertheorie abgeleitet. Der Gruppensatz erweist sich so als der einfachere Satz; in den Anwendungen des Gruppensatzes --- z.\ B. Bettische und Torsionszahlen in der Topologie --- ist somit ein Zurückgehen auf die Elementarteilertheorie nicht erforderlich.

\end{document}
